\section{Decoding as gradient separation}
\label{sec:decode-iff}

This section proves that, at a constrained stationary point on the LayerNorm sphere, the
single-token decoder selects the designated token if and only if a signed alignment
condition holds against every competitor. We first record an elementary sign equivalence,
then the main biconditional, then the descent-direction corollary for the
constrained-minimum case. Each displayed step carries a provenance comment
\verb|% @lx-from:| identifying the exact Lean declaration that verifies it.

% ===========================================================================
% Lemma: sign equivalence
% ===========================================================================
\begin{lemma}[Sign equivalence]
\label{lem:sign_equiv}
Let $\mu \in \mathbb{R}$ with $\mu \ne 0$, and let $c \in \mathbb{R}$ be arbitrary. Then % lint: ignore R15
\begin{align*}
\mu^{-1} c > 0 \;\iff\; \mu c > 0.
\end{align*}
\end{lemma}

\begin{proof}
% @lx-from: lem:sign_equiv | kind: arith
The two quantities differ by the strictly positive factor $\mu^{2}$: since $\mu \ne 0$ we
have $\mu^{2} > 0$, and
\begin{align*}
(\mu^{-1} c)\,\mu^{2} = \mu c.
\end{align*}
Multiplying a real number by the positive constant $\mu^{2}$ preserves its sign, so
$\mu^{-1} c$ and $\mu c$ are simultaneously positive, which is the claimed equivalence. This
completes the proof.
\end{proof}

% ===========================================================================
% Main theorem: decode iff gradient separation
% ===========================================================================
% @lx-from: thm:main | kind: named
\begin{theorem}[Decoding $\iff$ gradient separation]
\label{thm:main}
Let $\xstar \in \Sphere$ (so $\norm{\xstar} = r$ with $r > 0$) be constrained-stationary for
the loss with multiplier $\mu$, in the sense of \Cref{def:stationarity}:
\begin{align*}
\grad \loss(\xstar) = \mu\,\xstar,
\end{align*}
and suppose the multiplier is nonzero, $\mu \ne 0$. Then the designated token $\astar$ is
generated at $\xstar$ if and only if, for every competitor $b \ne \astar$, the signed
alignment of the gap row with the gradient is strictly positive:
\begin{align*}
\Generated(\astar, \xstar)
\;\iff\;
\forall\, b \ne \astar,\quad \mu\,\inner{W_{\astar} - W_{b}}{\grad \loss(\xstar)} > 0.
\end{align*}
\end{theorem}

\begin{proof}
The argument reduces the biconditional to a single competitor and is then quantified over
$b \ne \astar$. Throughout, write $\grad\loss(\xstar)$ for the (fixed) gradient vector; the
loss is never differentiated.

% @lx-from: thm:main.1 | kind: arith
First, since $\mu \ne 0$, the stationarity equation $\grad\loss(\xstar) = \mu\,\xstar$
rearranges to express $\xstar$ through the gradient:
\begin{align*}
\xstar = \mu^{-1}\,\grad \loss(\xstar).
\end{align*}

% @lx-from: thm:main.2 | kind: arith
Next, fix a competitor $b \ne \astar$. We claim the per-competitor cone inequality is
equivalent to the signed alignment:
\begin{align*}
\inner{W_{\astar}}{\xstar} > \inner{W_{b}}{\xstar}
\;\iff\;
\mu\,\inner{W_{\astar} - W_{b}}{\grad \loss(\xstar)} > 0.
\end{align*}
Indeed, by bilinearity the cone inequality $\inner{W_{\astar}}{\xstar} > \inner{W_{b}}{\xstar}$
is equivalent to $\inner{W_{\astar} - W_{b}}{\xstar} > 0$; substituting
$\xstar = \mu^{-1}\,\grad\loss(\xstar)$ and pulling the scalar out of the inner product gives
$\inner{W_{\astar} - W_{b}}{\xstar} = \mu^{-1}\,\inner{W_{\astar} - W_{b}}{\grad\loss(\xstar)}$,
so the inequality reads $\mu^{-1}\,\inner{W_{\astar} - W_{b}}{\grad\loss(\xstar)} > 0$.
Applying the sign equivalence of \Cref{lem:sign_equiv} with
$c = \inner{W_{\astar} - W_{b}}{\grad\loss(\xstar)}$ converts this to % lint: ignore R15
$\mu\,\inner{W_{\astar} - W_{b}}{\grad\loss(\xstar)} > 0$, as claimed.

% @lx-from: thm:main.3 | kind: named
Finally, unfolding the decoder predicate of \Cref{def:decoder},
$\Generated(\astar, \xstar)$ means exactly $\inner{W_{\astar}}{\xstar} > \inner{W_{b}}{\xstar}$
for all $b \ne \astar$. Quantifying the per-competitor equivalence just established over all
$b \ne \astar$ yields
\begin{align*}
\Generated(\astar, \xstar)
\;\iff\;
\forall\, b \ne \astar,\quad \mu\,\inner{W_{\astar} - W_{b}}{\grad \loss(\xstar)} > 0,
\end{align*}
which is the assertion. This completes the proof.
\end{proof}

\begin{remark}[Only $\mu \ne 0$ and stationarity are load-bearing]
\label{rem:generalization}
The proof of \Cref{thm:main} uses only the stationarity equation
$\grad\loss(\xstar) = \mu\,\xstar$ and the hypothesis $\mu \ne 0$ (which already forces
$\xstar \ne 0$). The sphere membership $\norm{\xstar} = r$ and the positivity $r > 0$ are part
of the fixed LayerNorm-sphere setting but are never invoked in the argument; the
characterization therefore holds verbatim for \emph{any} point satisfying the stationarity
equation with a nonzero multiplier, on or off the sphere. We retain $\norm{\xstar} = r$ and
$r > 0$ in the statement to match the intended geometric setting, and record their
non-use here as an honest generalization, not a defect.
\end{remark}

% ===========================================================================
% Corollary: descent-direction separation (mu < 0)
% ===========================================================================
% @lx-from: cor:neg | kind: named
\begin{corollary}[Descent-direction separation, $\mu < 0$]
\label{cor:neg}
Retain the hypotheses of \Cref{thm:main}, but assume the multiplier is negative, $\mu < 0$
(the constrained-minimum case). Then $\astar$ is generated at $\xstar$ if and only if the
descent direction $-\grad\loss(\xstar)$ is strictly better aligned with the answer row than
with every competitor row:
\begin{align*}
\Generated(\astar, \xstar)
\;\iff\;
\forall\, b \ne \astar,\quad \inner{W_{\astar} - W_{b}}{-\grad \loss(\xstar)} > 0.
\end{align*}
\end{corollary}

\begin{proof}
% @lx-from: cor:neg.1 | kind: named
Since $\mu < 0$, in particular $\mu \ne 0$, so \Cref{thm:main} applies with this $\xstar$ and
$\mu$, giving
\begin{align*}
\Generated(\astar, \xstar)
\;\iff\;
\forall\, b \ne \astar,\quad \mu\,\inner{W_{\astar} - W_{b}}{\grad \loss(\xstar)} > 0.
\end{align*}

% @lx-from: cor:neg.2 | kind: arith
It remains to rewrite the per-competitor condition. Fix $b \ne \astar$ and write
$v = W_{\astar} - W_{b}$. We show
\begin{align*}
\mu\,\inner{v}{\grad \loss(\xstar)} > 0
\;\iff\;
\inner{v}{-\grad \loss(\xstar)} > 0.
\end{align*}
Because $\mu < 0$, the product $\mu\,\inner{v}{\grad\loss(\xstar)}$ is strictly positive
exactly when $\inner{v}{\grad\loss(\xstar)} < 0$; and by linearity of the inner product in its
second argument, $\inner{v}{-\grad\loss(\xstar)} = -\inner{v}{\grad\loss(\xstar)}$, which is
strictly positive under precisely the same condition $\inner{v}{\grad\loss(\xstar)} < 0$.
Quantifying this equivalence over all $b \ne \astar$ turns the conclusion of \Cref{thm:main}
into the descent-direction form, which is the assertion. This completes the proof.
\end{proof}
